\section{Introduction}
\label{sec:introduction}

Financial markets are often studied as complex systems in which heterogeneous agents, institutional constraints, liquidity cycles, information flows and feedback mechanisms interact across multiple time scales \cite{mantegna_stanley_1999_book,cont_2001_empirical}. Even when the empirical object is reduced to daily asset returns, market variation should not be represented only by the marginal behaviour of each series. In practice, the variation of financial returns available in a market must also be described through the structure of relations among assets, which is summarized by their correlation matrix. This matrix reflects a mixture of collective market motion, sectoral organization, idiosyncratic shocks, crisis amplification and sampling noise. A central task of empirical econophysics is therefore to distinguish interpretable dependency structure from fluctuations that are compatible with random estimation effects \cite{raddant_dimatteo_2023}.

This distinction is relevant because financial risk is inherently multivariate. The volatility of an individual asset matters, but portfolio risk, contagion, diversification and systemic vulnerability depend on how assets co-move \cite{billio_2012_econometric}. During calm periods, correlations may appear moderate and dispersed. During stress periods, correlations often increase, reducing diversification benefits and lowering the effective dimensionality of the market \cite{junior_2012_correlation,zheng_2012_changes}. A purely univariate analysis cannot represent this collective component, whereas an unfiltered analysis of the full empirical correlation matrix may confuse persistent structure with noise.

Emerging equity markets are particularly relevant in this context. They combine global risk channels with local institutional, macroeconomic and liquidity conditions. Brazil provides a suitable empirical setting because B3 includes globally exposed commodity companies, large financial institutions, regulated utilities, consumer cyclicals, industrial firms, telecommunications firms and real-estate companies. These assets share exposure to domestic monetary policy, exchange-rate dynamics, country risk and global liquidity, but they also respond to distinct sectoral fundamentals. These characteristics suggest that the empirical dependency matrix may combine a dominant market component, group or sectoral components and a large residual component. Despite these empirical features, the Brazilian equity market remains relatively underexplored in the complex-systems literature. Previous studies have addressed Brazilian interbank networks \cite{tabak_2010_topology}, global financial system linkages \cite{silva_2016_structure}, agrarian cross-correlations \cite{stosic_2015_correlations}, multiscale networks \cite{pereira_2019_multiscale} and equity-market topology \cite{leonel_2022_brazilian}, but a broad B3 equity pipeline combining Random Matrix Theory filtering, planar financial networks and volatility forecasting remains insufficiently developed.

The empirical difficulty is that correlation matrices are noisy, especially when the number of assets is not negligible relative to the number of observations. Random Matrix Theory (RMT) provides a benchmark for this problem by comparing the empirical eigenspectrum with the spectrum expected from random correlation matrices \cite{marcenko_1967_distribution,laloux_1999_noise,plerou_2002_random}. Eigenvalues lying inside the Mar\v{c}enko--Pastur band are interpreted as compatible with random sampling noise under the benchmark model, whereas eigenvalues outside the band indicate candidate collective modes that deserve further interpretation. Subsequent work has extended RMT-based methods for cleaning large covariance matrices \cite{bun_2017_cleaning}. RMT is not used here as a mechanical classifier of economic sectors, but as an intermediate layer for reconstructing broad market, group-level and noise-compatible components before graph construction.

Network filtering provides a complementary representation of the same dependency system. By transforming correlations into distances, financial networks retain selected dependency links among assets and make it possible to inspect market backbones, local clusters and central nodes. The Minimum Spanning Tree (MST) emphasizes the sparsest connected structure \cite{mantegna_1999_hierarchical}, while the Planar Maximally Filtered Graph (PMFG) retains a denser planar representation with richer local connectivity \cite{tumminello_2005_tool,tumminello_2010_correlation}. These constructions allow one to ask whether strong dependencies are organized by sectors, whether central nodes change after market-mode removal and whether the mesoscopic topology of the market contains interpretable information beyond average correlation. Filtered networks are also relevant for risk analysis, where peripheral assets can offer better diversification properties than central hubs \cite{pozzi_2013_spread,aste_2010_correlation}.

This paper develops an integrated econophysics pipeline for Brazilian equities traded on B3. The analysis proceeds through three stages. The first stage documents stylized facts, static correlations, sectoral correlation differences and rolling market synchronization. The second stage applies RMT to decompose the empirical correlation matrix into market, group/sector and noise components, followed by hierarchical clustering and ordered heatmaps. The third stage converts original and filtered matrices into MST and PMFG networks, studies hub reorganization and aggregates dependencies at the subsector level. A volatility forecasting experiment then evaluates whether the extracted features provide incremental predictive information relative to established econometric benchmarks \cite{bollerslev_1986_generalized,corsi_2009_simple,hansen_lunde_2005}.


The contribution of this paper is fourfold:
\begin{enumerate}
  \item It builds a reproducible econophysics pipeline for B3 equities using adjusted daily prices from 1998 to 2025, including return construction, stylized-fact diagnostics, correlation analysis, sectoral comparison and rolling synchronization.
  \item It applies Random Matrix Theory to a 58-asset historical core universe and reconstructs market, group/sector and noise-compatible components of the empirical correlation matrix, providing a spectral interpretation of Brazilian equity dependencies.
  \item It compares MST and PMFG networks built from original and RMT-filtered matrices, with emphasis on hub reorganization, clique structure, clustering, sectoral retention and subsector-level aggregation.
  \item It evaluates whether RMT- and network-derived features provide incremental predictive information for realized-volatility forecasting relative to GARCH(1,1), HAR-RV, EWMA and classical volatility predictors, using a strict temporal split and nested feature-set design.
\end{enumerate}

The article is organized as follows. Section~\ref{sec:related_work} reviews financial dependencies, Random Matrix Theory, filtered financial networks and volatility forecasting. Section~\ref{sec:data} describes the data and asset universe. Section~\ref{sec:methodology} presents the empirical methodology and the mathematical definitions used throughout the paper. Sections~\ref{sec:stylized_correlation}--\ref{sec:forecasting} report the empirical results. Sections~\ref{sec:discussion} and~\ref{sec:limitations} discuss interpretation and limitations, and Section~\ref{sec:conclusion} presents the final considerations.
